\newpage
\section[Estado del arte]{Estado del Arte}

El desarrollo de un asistente educativo bilingüe español-Quechua Collao requiere la integración de avances en múltiples dominios disciplinares. A continuación, se presenta una síntesis analítica y crítica de las investigaciones previas más relevantes, organizadas temáticamente. Las fichas de análisis detalladas de cada uno de estos trabajos se presentan en el \hyperref[sec:anexos]{Anexo: Fichas de Análisis}.

\subsection{Procesamiento de Lenguaje Natural y Traducción Automática para Quechua Collao}
El procesamiento de lenguaje natural (PLN) para lenguas aglutinantes de bajos recursos como el quechua ha presentado desafíos históricos debido a su complejidad morfológica. Investigaciones recientes han abordado estas brechas mediante la construcción de corpus y el ajuste de modelos neuronales. \citet{Orcotoma2025} aportaron significativamente al construir un corpus paralelo Quechua Collao-español de 44\,263 palabras mediante el uso de reconocimiento óptico de caracteres (OCR) sobre textos académicos de la Academia Mayor de la Lengua Quechua, logrando tasas de error de palabra (WER) de 6.59\%. Por su parte, \citet{Duran2026} propusieron métodos semi-automáticos para la construcción de diccionarios quechua-español utilizando formalismos lingüísticos, lo que amplía la disponibilidad de recursos léxicos estructurados.

En cuanto a la traducción, \citet{SullaTorres2024} demostraron que los modelos de traducción automática neuronal (NMT) optimizados con BERT pueden alcanzar puntajes aceptables (BLEU de 20.13) en la dirección quechua-español, aunque subrayaron las limitaciones impuestas por la escasez de datos. No obstante, el paradigma de traducción y PLN ha cambiado con la llegada de los Grandes Modelos de Lenguaje (LLM). En este contexto, \citet{Court2024} evaluaron las habilidades de aprendizaje en contexto de LLMs comerciales (GPT-4, Gemini) para traducir del quechua al español. Sus hallazgos indican que, si bien la inclusión de traducciones a nivel de morfema mejora el rendimiento, la inyección de largas reglas gramaticales o ejemplos de corpus mediante técnicas de recuperación puede, paradójicamente, degradar el rendimiento de los modelos más avanzados, generando traducciones fluidas pero infieles o estereotipadas. Esto evidencia que depender exclusivamente de LLMs en la nube mediante \textit{prompting} no resuelve el problema de comprensión estructural en lenguas de bajos recursos.

\subsection{Small Language Models y Despliegue Offline}
Frente a la dependencia de conectividad y alto consumo computacional de los LLMs, el desarrollo de Small Language Models (SLM) y técnicas de optimización para hardware restringido (on-device AI) ha cobrado relevancia. La exhaustiva revisión de \citet{Xu2024} sistematiza cómo técnicas como la compartición de parámetros, la destilación de conocimiento y, fundamentalmente, la cuantización (PTQ y QAT), permiten que modelos con menos de 10 mil millones de parámetros se ejecuten eficientemente en dispositivos móviles y procesadores de bajo consumo.

La aplicación de estas tecnologías en contextos de baja conectividad fue validada empíricamente por \citet{Walusimbi2026}, quienes diseñaron \textit{Arapai}, una arquitectura de chatbot educativo orientado al funcionamiento offline. Al desplegar modelos cuantizados (como TinyLlama o Qwen) en hardware heredado (solo CPU), lograron latencias de respuesta de 1 a 3 segundos, demostrando que la inferencia local es técnicamente factible y pedagógicamente útil para el aprendizaje autodirigido en entornos con infraestructura precaria. Sin embargo, su estudio se centró en inglés y no abordó los retos de lenguas aglutinantes.

\subsection{Generación Aumentada por Recuperación (RAG) y Tecnología Educativa con IA}
La Generación Aumentada por Recuperación (RAG) se ha convertido en el estándar para mitigar las alucinaciones en modelos generativos, vinculando sus respuestas a corpus documentales específicos. La revisión sistemática de \citet{Li2025} confirma que los sistemas RAG mejoran significativamente la precisión factual y la personalización en aplicaciones educativas, permitiendo que la IA asuma roles de tutoría basados en el currículo local. 

Sin embargo, el despliegue de RAG con modelos pequeños presenta desafíos particulares. \citet{Hevia2025} evaluaron un tutor de IA offline combinando RAG con SLMs (SmolLM), encontrando que los modelos pequeños tienen dificultades para procesar contextos recuperados extensos o ruidosos, lo que irónicamente redujo su precisión en evaluaciones de respuestas múltiples. Esto sugiere la necesidad de estrategias de recuperación y segmentación documental altamente optimizadas. Adicionalmente, el despliegue de estas herramientas con estudiantes requiere garantías pedagógicas. \citet{Levonian2025} establecen que los sistemas de tutoría inteligente deben ser diseñados para generar interacciones seguras y relevantes, priorizando el andamiaje cognitivo sobre la simple provisión de respuestas directas.

\subsection{RAG vs.\ Fine-Tuning y Enfoque Híbrido para Lenguas de Bajos Recursos}
La decisión entre RAG y \textit{fine-tuning} como estrategia de inyección de conocimiento es central para esta tesis. \citet{Ovadia2024} encontraron que, para incorporar conocimiento factual nuevo, RAG superó consistentemente al \textit{fine-tuning} no supervisado, tanto para conocimiento existente como nuevo. Sin embargo, la evidencia no avala un RAG no optimizado: el cuello de botella de la tokenización en lenguas aglutinantes como el quechua infla el contexto requerido y penaliza la latencia, lo que exige un RAG disciplinado con fragmentos cortos y \textit{top-k} pequeño.

El enfoque RAFT (\textit{Retrieval-Augmented Fine-Tuning}) propuesto por \citet{Zhang2024RAFT} ofrece una síntesis prometedora: en lugar de entrenar al modelo a memorizar conocimiento, le enseña a usar el contexto RAG de forma resiliente, distinguiendo fragmentos relevantes de distractores. El antecedente más directo para esta tesis es el modelo QueEn \citep{Chen2024}, que demostró empíricamente que la hibridación de RAG sobre recursos lingüísticos combinada con adaptación LoRA elevó el puntaje BLEU de 1,5 a 17,6 en traducciones quechua--inglés, validando que la combinación de RAG con adaptación paramétrica ligera produce mejoras significativas para lenguas de bajos recursos.

\subsection{Política Educativa y Educación Intercultural Bilingüe en el Perú}
El desarrollo tecnológico debe alinearse con los marcos normativos y pedagógicos vigentes. El Modelo de Servicio Educativo Intercultural Bilingüe (MSEIB) del Ministerio de Educación del Perú \citep{minedu2018} establece lineamientos para garantizar una educación pertinente a estudiantes de pueblos originarios. El MSEIB distingue escenarios de fortalecimiento lingüístico (donde el quechua es lengua materna y el español segunda lengua) y revitalización (donde el español es lengua materna). Esto exige que cualquier herramienta de asistencia educativa no sea un simple traductor bidireccional, sino un sistema capaz de alternar códigos, detectar el idioma predominante del estudiante y recuperar contenidos con pertinencia cultural e intercultural, apoyando directamente la labor del docente bilingüe.

\subsection{Síntesis crítica y vacíos de investigación}
La revisión de literatura demuestra que existen avances maduros, pero aislados, en los componentes necesarios para el sistema propuesto. Por un lado, se han construido corpus y evaluado LLMs para quechua (\citealt{Orcotoma2025, Court2024}), y se ha demostrado que la hibridación de RAG con adaptación ligera produce mejoras significativas \citep{Chen2024}. Por otro lado, la viabilidad de SLMs \textit{offline} y sistemas RAG ha sido probada (\citealt{Walusimbi2026, Hevia2025, Xu2024}), pero principalmente para inglés y lenguas mayoritarias. La evidencia comparada confirma que RAG supera al \textit{fine-tuning} para inyección de conocimiento \citep{Ovadia2024}, pero que un RAG no optimizado con modelos pequeños puede degradar el rendimiento. La técnica RAFT \citep{Zhang2024RAFT} ofrece una vía para resolver este problema, pero no ha sido aplicada al Quechua Collao. Finalmente, la política EIB (\citealt{minedu2018}) exige recursos pertinentes que la actual oferta de LLMs comerciales en la nube no puede satisfacer, ya sea por barreras de conectividad o por la degradación del rendimiento al inyectarles reglas complejas \citep{Court2024}.

En consecuencia, el vacío de investigación se ubica en la ausencia de un prototipo integrado de asistente educativo offline español--Quechua Collao que combine un modelo compacto cuantizado, RAG local disciplinado, recursos lingüísticos validados y criterios pedagógicos EIB para apoyar consultas curriculares en contextos rurales del sur andino peruano.


